\documentclass[aspectratio=169,11pt]{beamer}
\usepackage[utf8]{inputenc}
\usepackage[T1]{fontenc}
\usepackage[spanish,es-noshorthands]{babel}
\usepackage{lmodern}
\usepackage{booktabs}
\usepackage{array}
\usepackage{listings}
\usetheme{metropolis}
\setbeamertemplate{navigation symbols}{}
\lstset{basicstyle=\ttfamily\scriptsize,columns=fullflexible,frame=single,breaklines=true}
\newcommand{\rel}[1]{\textbf{#1}}
\newcommand{\code}[1]{\texttt{#1}}

\title{Listwise Learning to Rank con AdaRank}
\subtitle{Training set de relevancia para consultas LOINC}
\author{[Integrante 1] \and [Integrante 2] \and [Integrante 3]}
\institute{Machine Learned Ranking --- MUIA, UPM}
\date{3 de diciembre de 2026}

\begin{document}

\maketitle

\begin{frame}{El problema}
  \begin{itemize}
    \item Tres consultas clínicas: \emph{glucose in blood}, \emph{bilirubin in plasma}, \emph{white blood cells count}.
    \item 67 pruebas de laboratorio LOINC candidatas, \emph{las mismas} para las tres consultas, lo que da 201 pares.
    \item Objetivo: ordenar los 67 documentos de cada consulta por relevancia.
    \item El Excel no trae etiquetas: hay que \textbf{construir el training set}.
  \end{itemize}
  \vspace{4pt}
  \small Cada documento tiene \code{loinc\_num}, \code{long\_common\_name}, \code{component}, \code{system} y \code{property}.
\end{frame}

\begin{frame}{Pointwise, pairwise, listwise}
  \begin{tabular}{>{\bfseries}l p{4.2cm} p{5cm}}
    \toprule
    Enfoque & Unidad de aprendizaje & Ejemplos \\
    \midrule
    Pointwise & un documento $\rightarrow$ etiqueta & regresión, clasificación ordinal \\
    Pairwise & par $(d_i,d_j)$: ¿cuál va antes? & RankNet, RankSVM \\
    Listwise & la lista completa de una consulta & \alert{AdaRank}, ListNet, LambdaMART \\
    \bottomrule
  \end{tabular}
  \vspace{8pt}
  \begin{itemize}
    \item Al usuario le importa el \emph{orden} de la lista, que es lo que optimiza un método listwise.
    \item Por eso el training set se agrupa por \code{qid}: 3 listas de 67, no 201 filas sueltas.
  \end{itemize}
\end{frame}

\begin{frame}{AdaRank (Xu y Li, 2007)}
  \begin{itemize}
    \item \emph{Boosting} sobre consultas: el ranker final es $f(\mathbf{x})=\sum_t \alpha_t h_t(\mathbf{x})$, y cada $h_t$ es una feature.
    \item Mantiene pesos $P_t(i)$ sobre las consultas. En cada ronda:
      \begin{enumerate}
        \item elige la feature con mejor NDCG ponderado por $P_t$;
        \item calcula $\alpha_t=\frac12\ln\frac{\sum_i P_t(i)(1+E_i)}{\sum_i P_t(i)(1-E_i)}$;
        \item actualiza $P_{t+1}(i)\propto e^{-E(q_i,f_t)}$, de modo que las consultas mal ordenadas pesan más.
      \end{enumerate}
    \item Optimiza directamente una medida de ranking:
      $\mathrm{NDCG@}k=\frac{1}{\mathrm{IDCG@}k}\sum_{r\le k}\frac{2^{y_r}-1}{\log_2(r+1)}$.
  \end{itemize}
\end{frame}

\begin{frame}{Escala de relevancia 0--3: tres dimensiones}
  \begin{itemize}
    \item \textbf{Analito} (\code{component}): exacto / relacionado (fracción o subpoblación) / ninguno.
    \item \textbf{Muestra} (\code{system}): exacta / compatible pero distinta / incompatible.
    \item \textbf{Medida} (\code{property}, sufijo \code{\^{}}): adecuada / desviada.
  \end{itemize}
  \vspace{6pt}
  \begin{block}{Regla determinista}
    Analito distinto $\rightarrow$ \rel{0}. Si no, base \rel{3} (exacto) o \rel{2} (relacionado);
    $-1$ si la muestra es compatible, $-2$ si es incompatible; $-1$ si la medida está desviada; mínimo \rel{1}.
  \end{block}
  \small Cada par tiene su justificación en \code{data/relevance\_labels.csv}.
\end{frame}

\begin{frame}{Decisiones de etiquetado}
  \begin{itemize}
    \item \code{Ser/Plas} frente a \code{Bld}: compatible pero distinta, así que la glucosa en suero/plasma es \rel{2}.
    \item \textbf{WBC = leucocitos}. Los recuentos de linfocitos y monocitos son \rel{2}; el \emph{porcentaje} de linfocitos es \rel{1}.
    \item La glucosa en ayunas o post-carga es una prueba funcional: medida desviada.
    \item ``Mismo término'' no implica ``misma prueba'':
      \begin{itemize}
        \item glucosa en orina o en líquido pleural $\rightarrow$ \rel{1};
        \item \emph{Blood product type} $\rightarrow$ \rel{0};
        \item \emph{Calcium bilirubinate in Stone} $\rightarrow$ \rel{0}.
      \end{itemize}
  \end{itemize}
\end{frame}

\begin{frame}{Etiquetas resultantes}
  \centering\scriptsize
  \resizebox{\linewidth}{!}{% Generado por src/export_tables.py — no editar a mano
\begin{tabular}{llp{6.2cm}llc}
\toprule
Consulta & LOINC & long\_common\_name & system & property & Rel. \\
\midrule
glucose in blood & 74774-1 & Glucose [Mass/volume] in Serum, Plasma or Blood & Ser/Plas/Bld & MCnc & 3 \\
 & 14749-6 & Glucose [Moles/volume] in Serum or Plasma & Ser/Plas & SCnc & 2 \\
 & 15076-3 & Glucose [Moles/volume] in Urine & Urine & SCnc & 1 \\
 & 14747-0 & Glucose [Moles/volume] in Pleural fluid & Plr fld & SCnc & 1 \\
 & 14764-5 & Glucose [Moles/volume] in Serum or Plasma --3 hours post 100 g glucose PO & Ser/Plas & SCnc & 1 \\
 & 35184-1 & Fasting glucose [Mass or Moles/volume] in Serum or Plasma & Ser/Plas & MSCnc & 1 \\
\midrule
bilirubin in plasma & 1975-2 & Bilirubin.total [Mass/volume] in Serum or Plasma & Ser/Plas & MCnc & 3 \\
 & 1968-7 & Bilirubin.direct [Mass/volume] in Serum or Plasma & Ser/Plas & MCnc & 2 \\
 & 1971-1 & Bilirubin.indirect [Mass/volume] in Serum or Plasma & Ser/Plas & MCnc & 2 \\
 & 35192-4 & Bilirubin.indirect [Mass or Moles/volume] in Serum or Plasma & Ser/Plas & MSCnc & 2 \\
 & 14423-8 & Bilirubin.total [Mass/volume] in Synovial fluid & Synv fld & MCnc & 1 \\
 & 33870-7 & Bilirubin.total [Presence] in Unspecified specimen & XXX & PrThr & 1 \\
\midrule
white blood cells count & 26464-8 & Leukocytes [\#/volume] in Blood & Bld & NCnc & 3 \\
 & 26474-7 & Lymphocytes [\#/volume] in Blood & Bld & NCnc & 2 \\
 & 26484-6 & Monocytes [\#/volume] in Blood & Bld & NCnc & 2 \\
 & 26478-8 & Lymphocytes/100 leukocytes in Blood & Bld & NFr & 1 \\
\bottomrule
\end{tabular}
}
\end{frame}

\begin{frame}{Ampliación sistemática: por qué}
  \begin{itemize}
    \item Una primera ampliación con términos buscados a mano tenía \textbf{sesgo de selección}:
      \begin{itemize}
        \item cada consulta tenía una lista construida de forma distinta;
        \item casi todos los candidatos contenían la palabra de la consulta.
      \end{itemize}
    \item ``Equilibrada'' \emph{no} significa repartir las notas por igual. Eso no es realista y distorsiona el NDCG.
    \item Lo correcto es aplicar \textbf{el mismo procedimiento a todas las consultas} y dejar que la proporción de relevantes salga sola.
  \end{itemize}
\end{frame}

\begin{frame}{Ampliación sistemática: 12 consultas estratificadas}
  \begin{itemize}\small
    \item Para cada estrato \textbf{dominio clínico $\times$ muestra}, la prueba más frecuente de LOINC (\code{COMMON\_TEST\_RANK}).
      El texto de la consulta se redacta a mano.
  \end{itemize}
  \centering\scriptsize
  \begin{tabular}{lll}
    \toprule
    Dominio & Muestra & Consulta \\
    \midrule
    Química & suero / sangre / orina & creatinine in serum; oxygen partial pressure in blood; ketones in urine \\
    Hematología & sangre & red blood cells count \\
    Coagulación & plasma & prothrombin time in plasma \\
    Microbiología & respiratoria / orina / suero & SARS-CoV-2 RNA; bacteria in urine culture; HIV Ab+Ag \\
    Uroanálisis & orina & leukocyte esterase in urine \\
    Serología & suero & perinuclear ANCA in serum \\
    Fármacos & sangre & tacrolimus in blood \\
    Patología & tejido & EGFR antigen in tissue \\
    \bottomrule
  \end{tabular}
\end{frame}

\begin{frame}{Ampliación sistemática: pooling y etiquetas}
  \begin{itemize}\small
    \item \textbf{Pooling} (como TREC/LETOR) sobre 62\,148 pruebas de laboratorio, igual para las 15 consultas:
      40 por BM25 + 20 por componente + 300 aleatorias, con límite de casi-duplicados. Resultado: \textbf{1222 documentos}.
    \item Misma regla 0--3; las 201 etiquetas oficiales no cambian. Salen distractores difíciles (rel. 0):
  \end{itemize}
  \centering\tiny
  \resizebox{!}{0.55\textheight}{% Generado por src/export_tables.py — no editar a mano
\begin{tabular}{lp{6.6cm}lllc}
\toprule
Consulta & Documento de la colección & Analito & Muestra & Medida & Rel. \\
\midrule
oxygen partial pressure in blood & Oxygen [Partial pressure] in Blood & exacto & exacta & adecuada & 3 \\
 & Oxygen saturation in Arterial blood & relacionado & exacta & desviada & 1 \\
\midrule
prothrombin time in plasma & Prothrombin time (PT) & exacto & exacta & adecuada & 3 \\
 & INR in Platelet poor plasma by Coagulation assay & relacionado & exacta & adecuada & 2 \\
 & aPTT in Blood by Coagulation assay & ninguno & compatible & adecuada & 0 \\
\midrule
ketones in urine & Ketones [Presence] in Urine & exacto & exacta & adecuada & 3 \\
 & Methyl ethyl ketone [Mass/volume] in Urine & ninguno & exacta & adecuada & 0 \\
\midrule
SARS-CoV-2 RNA in respiratory specimen & SARS-CoV-2 (COVID-19) RNA [Presence] in Respiratory system specimen & exacto & exacta & adecuada & 3 \\
 & SARS-CoV+SARS-CoV-2 (COVID-19) Ag [Presence] in Respiratory system specimen by Rapid immunoassay & relacionado & exacta & adecuada & 2 \\
 & SARS coronavirus RNA [Presence] in Specimen by NAA with probe detection & ninguno & incompatible & adecuada & 0 \\
\midrule
bacteria in urine culture & Bacteria identified in Urine by Culture & exacto & exacta & adecuada & 3 \\
 & Bacteria identified in Bronchoalveolar lavage by Culture & exacto & incompatible & adecuada & 1 \\
\midrule
HIV antibody and antigen in serum & HIV 1+2 Ab+HIV1 p24 Ag [Presence] in Serum or Plasma by Immunoassay & exacto & exacta & adecuada & 3 \\
 & HIV 1+2 Ab [Presence] in Serum or Plasma by Immunoassay & relacionado & exacta & adecuada & 2 \\
\midrule
tacrolimus in blood & Tacrolimus [Mass/volume] in Blood & exacto & exacta & adecuada & 3 \\
 & Tacrolimus [Mass/volume] in Blood --trough & exacto & exacta & desviada & 2 \\
\midrule
EGFR antigen in tissue & Epidermal growth factor receptor Ag [Presence] in Tissue & exacto & exacta & adecuada & 3 \\
 & Epidermal growth factor [Mass/volume] in Serum or Plasma & ninguno & incompatible & desviada & 0 \\
\bottomrule
\end{tabular}
}
\end{frame}

\begin{frame}[fragile]{Features (sin fuga de etiquetas)}
  \begin{tabular}{>{\bfseries}l l}
    \toprule
    $f_1$ & coseno TF-IDF consulta--\code{long\_common\_name} \\
    $f_2$ & BM25 normalizado por consulta \\
    $f_3$ & solapamiento con sinónimos (\emph{white cells}$\to$\emph{leukocytes}) \\
    $f_4$ & analito de la consulta en \code{component} (0/1) \\
    $f_5$ & término de muestra en \code{system} (0/1) \\
    \bottomrule
  \end{tabular}
  \vspace{6pt}
  \begin{itemize}\small
    \item Se eliminaron las codificaciones \code{LabelEncoder}: son índices alfabéticos que no dependen de la consulta.
    \item LETOR: \code{train.txt} (3$\times$67), \code{train\_terms.txt} (3$\times$1222), \code{train\_extended.txt} (15$\times$1222).
  \end{itemize}
  \vspace{-4pt}
  \begin{lstlisting}
3 qid:1 1:0.737931 2:1 3:1 4:1 5:1 # docid = 74774-1
2 qid:1 1:0.537084 2:0.694848 3:0.5 4:1 5:0 # docid = 14749-6
  \end{lstlisting}
\end{frame}

\begin{frame}{Ranking del modelo final (3 consultas)}
  \centering\scriptsize
  \resizebox{!}{0.8\textheight}{% Generado por src/export_tables.py — no editar a mano
\begin{tabular}{lrp{7cm}cr}
\toprule
Consulta & Pos. & Documento & Rel. & Score \\
\midrule
glucose in blood & 1 & Glucose [Mass/volume] in Serum, Plasma or Blood & 3 & 3.49 \\
 & 2 & Glucose [Moles/volume] in Serum or Plasma & 2 & 2.22 \\
 & 3 & Glucose [Moles/volume] in Urine & 1 & 1.91 \\
 & 4 & Glucose [Moles/volume] in Serum or Plasma --3 hours post 100 g glucose PO & 1 & 1.84 \\
 & 5 & Fasting glucose [Mass or Moles/volume] in Serum or Plasma & 1 & 1.77 \\
 & 6 & Glucose [Moles/volume] in Pleural fluid & 1 & 1.71 \\
\midrule
bilirubin in plasma & 1 & Bilirubin.total [Mass/volume] in Serum or Plasma & 3 & 3.18 \\
 & 2 & Bilirubin.indirect [Mass/volume] in Serum or Plasma & 2 & 3.09 \\
 & 3 & Bilirubin.direct [Mass/volume] in Serum or Plasma & 2 & 2.97 \\
 & 4 & Bilirubin.indirect [Mass or Moles/volume] in Serum or Plasma & 2 & 2.97 \\
 & 5 & Bilirubin.total [Mass/volume] in Synovial fluid & 1 & 1.70 \\
 & 6 & Bilirubin.total [Presence] in Unspecified specimen & 1 & 1.69 \\
\midrule
white blood cells count & 1 & Leukocytes [\#/volume] in Blood & 3 & 2.71 \\
 & 2 & Lymphocytes [\#/volume] in Blood & 2 & 2.00 \\
 & 3 & Monocytes [\#/volume] in Blood & 2 & 1.92 \\
 & 4 & Lymphocytes/100 leukocytes in Blood & 1 & 1.91 \\
 & 5 & ABO group [Type] in Blood from Blood product unit & 0 & 1.73 \\
 & 6 & Blood product type & 0 & 1.54 \\
\bottomrule
\end{tabular}
}
\end{frame}

\begin{frame}{Evaluación leave-one-query-out}
  \centering\scriptsize
  \resizebox{\linewidth}{!}{% Generado por src/export_tables.py — no editar a mano
\begin{tabular}{lrrrr}
\toprule
Sistema & NDCG@1 & NDCG@3 & NDCG@5 & NDCG@10 \\
\midrule
\multicolumn{5}{l}{\emph{3 consultas oficiales (3 folds)}} \\
\textbf{AdaRank} & 0.667 & 0.772 & 0.796 & 0.797 \\
solo $f_1$ (coseno TF-IDF) & 0.667 & 0.856 & 0.889 & 0.889 \\
solo $f_2$ (BM25) & 0.667 & 0.764 & 0.829 & 0.863 \\
solo $f_3$ (solap. sinónimos) & 1.000 & 0.875 & 0.885 & 0.877 \\
solo $f_4$ (analito en component) & 0.714 & 0.824 & 0.821 & 0.824 \\
solo $f_5$ (muestra en system) & 0.333 & 0.366 & 0.475 & 0.467 \\
\midrule
\multicolumn{5}{l}{\emph{6 consultas, ampliado$^\dagger$ (6 folds)}} \\
\textbf{AdaRank} & 0.833 & 0.928 & 0.944 & 0.944 \\
solo $f_1$ (coseno TF-IDF) & 0.833 & 0.928 & 0.944 & 0.944 \\
solo $f_2$ (BM25) & 0.833 & 0.882 & 0.914 & 0.931 \\
solo $f_3$ (solap. sinónimos) & 0.905 & 0.878 & 0.878 & 0.893 \\
solo $f_4$ (analito en component) & 0.619 & 0.820 & 0.832 & 0.834 \\
solo $f_5$ (muestra en system) & 0.167 & 0.183 & 0.237 & 0.234 \\
\bottomrule
\end{tabular}
}
  \vspace{4pt}
  \begin{itemize}\small
    \item Con 3 consultas, AdaRank sobreajusta a las 2 de entrenamiento y no supera a la mejor feature sola.
    \item Con 15 consultas variadas queda cerca de las mejores features (0.627 frente a 0.660 de $f_3$). Va bien si el nombre LOINC
      contiene la consulta (EGFR 1.00) y falla si no (\emph{red blood cells count}: 0.00, porque \emph{MCHC in Red Blood Cells} gana).
  \end{itemize}
\end{frame}

\begin{frame}{El pooling genera distractores difíciles}
  \centering\scriptsize
  % Generado por src/export_tables.py — no editar a mano
\begin{tabular}{lrp{7cm}cr}
\toprule
Consulta & Pos. & Documento & Rel. & Score \\
\midrule
white blood cells count & 1 & Plasma cells/Leukocytes in Blood by Manual count & 1 & 0.88 \\
 & 2 & Other cells/Leukocytes in Blood by Manual count & 1 & 0.81 \\
 & 3 & Other cells/Leukocytes in Blood by Automated count & 1 & 0.81 \\
 & 4 & Mononuclear cells/Leukocytes in Blood by Manual count & 1 & 0.81 \\
 & 5 & Malignant cells/Leukocytes in Blood by Manual count & 1 & 0.80 \\
 & 6 & Smudge cells/Leukocytes in Blood by Manual count & 1 & 0.79 \\
\multicolumn{5}{c}{$\vdots$} \\
 & 21 & Leukocytes [\#/volume] in Blood by Automated count & 3 & 0.75 \\
\bottomrule
\end{tabular}

  \vspace{6pt}
  \begin{itemize}\small
    \item \emph{Plasma cells/Leukocytes in Blood} contiene literalmente \emph{cells} y \emph{blood}; el recuento total de leucocitos, no.
    \item Las features léxicas no distinguen los matices que separan un 3 de un 1. Hacen falta features más expresivas.
  \end{itemize}
\end{frame}

\begin{frame}{Limitaciones y conclusiones}
  \begin{itemize}
    \item Con 3--15 consultas los resultados son demostrativos; repetir documentos no crea consultas nuevas.
    \item Hacen falta más consultas \emph{y} features más expresivas (estructurales, semánticas) para los términos difíciles.
    \item $f_4$ y $f_5$ se parecen a la regla de etiquetado, aunque son más gruesas.
    \item Corregimos un bug de AdaRank: $\alpha=\infty$ cuando un ranker débil es perfecto.
    \item[$\Rightarrow$] Entregamos un training set listwise reproducible: regla 0--3 documentada, features, LETOR,
      validaciones y AdaRank entrenado y evaluado.
  \end{itemize}
\end{frame}

\begin{frame}[standout]
  ¿Preguntas?
\end{frame}

\end{document}
